\section{Introduction}
\label{sec:introduction}

Recent progress in artificial intelligence (AI) has resulted in rapidly improved AI capabilities. These capabilities are not designed by humans. Instead, they are learned by deep neural networks \citep{Hinton_2006_EarlyDNN, LeCun_2015_DeepLearning}. Developers only need to design the training process; they do not need to -- and in almost all cases, do not -- understand the neural mechanisms underlying the capabilities learned by an AI system.

Although human understanding of these mechanisms is not necessary for AI capabilities, understanding them would enhance several \textit{human} abilities. For example, it would permit better human control over AI behavior and better monitoring during deployment. It would also facilitate trust in AI systems, allowing us to fully realize their potential benefits by enabling their deployment in safety-critical and ethically-sensitive settings.

Beyond the engineering benefits, understanding AI systems offers immense scientific opportunities. For the first time in history, we can create and study artificial minds with a level of access and control that is simply not possible in biological systems. What new laws of nature governing the mechanisms of minds might we discover from studying the internal workings of AI systems?

The scientific opportunities are not limited to the field of AI. If an AI can outperform tools designed by humans in a given scientific field, it suggests that the AI system is representing something about the world currently unknown to us. We can develop a deeper comprehension of the world by understanding those representations. What can we learn about protein folding from AIs that can successfully predict protein structure? What insights can we glean about disease from a radiographer that performs beyond human ability?

\textit{Mechanistic interpretability} might unlock these benefits. This field of study aims to understand neural networks' decision-making processes. Here, we define “Understanding a neural network's decision-making process” as the ability to use knowledge about the mechanisms underlying a network's decision-making process in order to successfully predict its behavior (even on arbitrary inputs) or to accomplish other practical goals with respect to the network. Such goals might include more precise control of the network's behavior, or improved network design. Interpretability promises greater assurances for AI systems through a better understanding of what neural networks have learned, thus enabling us to realize their potential benefits.

\subsection{The focus of this review: Open problems and the future of mechanistic interpretability}
\label{subsec:review-focus}

Several recent reviews of mechanistic interpretability research and related topics exist \citep{rauker2023transparentaisurveyinterpreting,Geiger_2022_CausalStructure,Bereska_2024_MechInterpReview, ferrando2024primerinnerworkingstransformerbased,rai2024practicalreviewmechanisticinterpretability,anwar2024foundationalchallengesassuringalignment,davies2024cognitiverevolutioninterpretabilityexplaining,mosbach2024insightsactionsimpactinterpretability, mueller2024questrightmediatorhistory}. Our review takes a more forward-looking stance. We discuss not only where the frontier is today, but also which directions we might benefit most from prioritizing in the future.

\subsubsection{Why `mechanistic' interpretability?}
\label{subsec:mechanistic}

The distinction between interpretability and mechanistic interpretability is not always clear and is therefore worth clarifying.
The motivations and methods used in interpretability work are often diverse \citep{Lipton_2016_InterpretabilityMotivations,doshivelez2017rigorousscienceinterpretablemachine,Jacovi_2023_ExplainableAiLiterature}. As a result, there are many ways in which interpretability research might be categorized. Prior categorizations of interpretability include causal vs. correlational methods, supervised vs. unsupervised methods, bottom-up vs. top-down methods, among others \citep{geiger2021causalabstractionsneuralnetworks, mueller2024questrightmediatorhistory,Bereska_2024_MechInterpReview,Belinkov_2022_ProbingClassifiers,zou2023representationengineeringtopdownapproach, davies2024cognitiverevolutioninterpretabilityexplaining}. This review focuses specifically on \textit{mechanistic} interpretability. But what distinguishes `mechanistic interpretability' from interpretability in general? It has been noted that the term is used in a number of (sometimes inconsistent) ways \citep{saphra2024mechanistic}. In this review, we use the term `mechanistic interpretability' in a technical sense, referring specifically to work that investigates the mechanisms underlying neural network generalization. Mechanistic interpretability represents one of three threads of interpretability research, each with distinct but sometimes overlapping motivations, which roughly reflects the changing aims of interpretability work over time.

The first thread aims to build AI systems that are interpretable by design. Much early interpretability work focused on explaining the sensitivity of machine learning models to inputs and training data. This work typically used small-to-medium sized models designed to be easily interpretable, such as decision trees \citep{Breiman_1984_Classification_Regression_Trees, Hu_2019_OptimalSparseDecisionTrees}, linear models \citep{Roweis_1999_LinearModel}, and generalized additive models \citep{Hastie_1986_AdditiveModel,Agarwal_2020_Additive_Model}. These models could be used alongside attribution methods such as influence functions \citep{Hampel_1986_Influence_Function,Koh_2017_influence_function} and Shapley values \citep{Shapley_1997_Shapley_Values, Lundberg_2017_SHAP}, which were common techniques used to characterize model decision boundaries with respect to inputs. “Interpretability by design” continues to be an active research area, including architectures such as Concept-Bottleneck Models \citep{Koh_2020_ConceptBottleneckModels}, Backpack Language Models \citep{Hewitt_2023_BackpackLanguageModels},
Kolmogorov-Arnold Networks \citep{Liu_2024_KolmogorovArtificialNetwork}, and sparse decision trees \citep{Xin_2022_SparseDecisionTrees}.

With the rise of larger-scale nonlinear neural networks \citep{Krizhevsky_2012_ImageNet, He_2015_ImageDNN}, another thread grew in importance, driven primarily by the question: Why did my model make this particular decision? However, one challenge in interpreting larger networks was finding attribution methods that could scale to large networks \citep{Zeiler_2013_VisualizingConvolutionalNetworks}. In response, a number of local attribution methods were developed, including grad-CAM \citep{Selvaraju_2019_GradCAM}, integrated gradients \citep{Sundarajan_2017_IntegratedGradients}, and masking-based causal attribution \citep{Fong_2017_MaskingBasedCausalAttribution}, SHAP \citep{Lundberg_2017_SHAP}, LIME \citep{Ribiero_2016_LIME}, and many other methods, including backprop-based visualization methods \citep{Simonyan_2014_BackpropVisualisation, Nguyen_2016_FeatureVisualization}.


Inspired by, for example, Inception \citep{szegedy2015inception} and GPT-3 \citep{Brown_2020_GPT3}, another thread emerged as models became capable of more profound generalization. Focused on the broader subject of generalization, it was driven by the question: How did my model solve this general class of problems? Due to its emphasis on the mechanisms underlying neural network generalization, the work in this category is commonly referred to as `mechanistic interpretability' (in addition to other, more cultural reasons, according to \citep{saphra2024mechanistic}). This kind of interpretability work is driven by a fundamental hypothesis in deep learning that generalization arises from shared computation \citep{LeCun_2015_DeepLearning}. Early work in this area, such as feature visualization \citep{olah2017feature}, or network dissection \citep{Bau_2017_NetworkDissection}, sought “global” explanations for model generalization by investigating the roles of model components across a class of decisions. More recent work in this area looks at “circuits” of components \citep{Wang_2022_Circuits}, generalizable patterns of information flow \citep{Geva_2023_InformationFlow}, representation subspaces
\citep{geiger2024findingalignmentsinterpretablecausal,zou2023representationengineeringtopdownapproach} and probes (or self-supervised searches via sparse dictionary learning) for representation vectors that carry information that generalizes across many instances for a particular task or set of tasks 
\citep{cunningham2023sparseautoencodershighlyinterpretable,Bricken_2023_dictionary,Todd_2024_Probe,Tigges_2023_Probe}.

\subsection{Types of open problems}
\label{subsec:types-of-open-problems}


The field of mechanistic interpretability ultimately aims to achieve concrete scientific and engineering goals. For instance, we would like to be able to:
\begin{itemize}
    \item Monitor AI systems for signs of cognition related to dangerous behavior (\Cref{subsec:monitoring-and-auditing});
    \item Modify internal mechanisms and edit model parameters to adapt their behavior to better suit our needs (\Cref{subsec:better-control});
    \item Predict how models will act in unseen situations or predict when a model might learn specific abilities (\Cref{subsec:better-predictions});
    \item Improve model inference, training, and mechanisms to better suit our preferences (\Cref{subsec:capabilities});
    \item Extract latent knowledge from AI systems so we can better model the world (\Cref{subsec:microscope-ai}).
\end{itemize}

Despite recent hopeful signs of progress, mechanistic interpretability still has considerable distance to cover before achieving satisfactory progress toward most of its scientific and engineering goals.

To achieve these goals, the field not only needs greater application of current state-of-the-art mechanistic interpretability methods, but also requires the development of improved techniques. The first major section (\Cref{sec:open-problems-methods-and-foundations}) therefore discusses open problems related to the methods and foundations of mechanistic interpretability.

We then explore key axes of research progress that will determine how far we can advance toward the goals of mechanistic interpretability (\Cref{subsec:axes}). \Cref{sec:open-problems-applications} outlines how applications of interpretability methods have made progress toward the field's goals, and, for each goal, discusses the specific axes along which progress is likely needed to achieve specific objectives.

Finally, we note that the goals, applications, and methods of mechanistic interpretability do not exist in a vacuum. Like any scientific field, they lie within a broader societal context. The final section of this review examines open socio-technical problems in mechanistic interpretability (\Cref{sec:open-problems-sociotechnical}). It discusses current initiatives and possible pathways to translate technical progress into levers for AI governance, alongside consequential social and philosophical challenges faced by the field.
